% Accepted after independent root and intervals re-proofs; integrated into sections/C.tex.
\begin{claim}{C.26}{Refuted}{A three-sort counterexample to the unrestricted linear theorem}\label{C26}
There is a complete and cocomplete $k$-linear abelian category $\mathcal M$, over any nonzero field $k$, with enough projectives and a cocontinuous symmetric $k$-bilinear tensor, satisfying all four hypotheses (a)--(d) at 134-8/048, for which the canonical functor
\[
\mathcal M[\mathcal W^{-1}]\longrightarrow D_{\geq0}(k),
\qquad \mathcal W=(L\varphi)^{-1}(\mathrm{qis}),
\]
is not an equivalence. This concerns the unrestricted theorem, not its corollary with finitely presented projective generators.
\end{claim}
\begin{proof}
\textbf{The abelian category and its scalar adjoints.}
Put $E(V)=V^{**}/V$ as in Claim~\ref{C24}. We use its proved exactness, $E(k)=0$, and $E(W)\ne0$ for an infinite-dimensional vector space $W$. On $\operatorname{Vect}_k^3$ take
\[
H(U_0,U_1,U_2)=(E(U_2),U_0,U_1),\qquad T=\id\oplus H,
\]
with square-zero multiplication on the $H$-summand. Its algebra category $\mathcal B$ has objects
\[
(U_0,U_1,U_2,a:U_0\to U_1,c:U_1\to U_2,b:E(U_2)\to U_0),
\quad ca=0,\ ab=0,\ bE(c)=0.
\]
Morphisms are triples commuting with these operations. Exactness of $E$ gives kernels and cokernels componentwise, so $\mathcal B$ is abelian and its forgetful functor to $\operatorname{Vect}_k^3$ is exact and faithful. All small limits are created by that functor. Products are exact, since the underlying products of vector-space triples are exact and the forgetful functor preserves and reflects exactness. The free functor has underlying functor $T$, is exact, and takes every vector-space triple to a projective object. Its counit is epimorphic, hence there are enough projectives. Arbitrary coproducts exist by the free coequalizer presentation: for each object use the two maps from the free algebra on $T$ of its underlying triple to the free algebra on its underlying triple, take free algebras on the direct sums for a family, and take the cokernel of their difference. This proves cocompleteness without asserting that underlying coproducts are componentwise. These are the same elementary arguments verified in Claim~\ref{C24}, with three sorts in place of two.

Set
\[
\kappa V=(V,0,0),\qquad
\phi(U_0,U_1,U_2,a,c,b)=\coker b,\qquad
\psi(U_0,U_1,U_2,a,c,b)=\ker a.
\]
The morphism equations give $\phi\dashv\kappa\dashv\psi$. In particular $\kappa$ is the actual scalar-copower functor of $1=\kappa k$, is fully faithful and exact, and preserves all small limits and colimits.

Write $S_i(V)$ for the triple with $V$ in sort $i$ only and zero operations. The free projectives on the three sorts are
\[
P_0(V)=(V,V,0;\ a=\id),\quad
P_1(V)=(0,V,V;\ c=\id),\quad
P_2(V)=(E(V),0,V;\ b=\id).
\]
Every projective is a retract of a direct sum of these. We abbreviate $S_1$ to $S$. There are exact sequences
\begin{align*}
0&\longrightarrow S(V)\xrightarrow{i_V}P_0(V)\longrightarrow\kappa V\longrightarrow0,\\
0&\longrightarrow S_2(V)\longrightarrow P_1(V)\xrightarrow{q_V}S(V)\longrightarrow0,\\
0&\longrightarrow\kappa E(V)\longrightarrow P_2(V)\longrightarrow S_2(V)\longrightarrow0.
\end{align*}
The functor $S$ preserves arbitrary coproducts: its presentation is the cokernel of the natural map $P_2(V)\to P_1(V)$, and both free functors preserve coproducts. This observation is needed in the tensor construction.

\textbf{The canonical strength.}
Every $k$-linear functor $E$ on vector spaces has a canonical natural map
\[
t_{V,W}:E(V)\otimes_k W\longrightarrow E(V\otimes_k W),
\quad t_{V,W}(e\otimes w)=E(v\mapsto v\otimes w)(e).
\]
Its bilinearity follows from $k$-linearity of $E$. Naturality, its unit identity at $W=k$, and
\[
t_{V\otimes W,Z}(t_{V,W}\otimes Z)=t_{V,W\otimes Z}
\]
follow by evaluating on $e\otimes w\otimes z$. There is a left-hand version obtained by the ordinary vector-space symmetry. We use only this strength, never a costrength in the reverse direction.

\textbf{A bifunctor on all free projectives.}
For $L=U\otimes_k V$ prescribe
\[
\begin{array}{c|ccc}
\boxtimes&P_0(V)&P_1(V)&P_2(V)\\ \hline
P_0(U)&P_0(L)\oplus S(L)&P_1(L)&P_2(L)\\
P_1(U)&P_1(L)&0&0\\
P_2(U)&P_2(L)&0&0.
\end{array}
\]
We specify the maps as well; the object table alone would not define a tensor. In addition to diagonal coefficient maps, every map between free projectives is a sum of maps of the following three types:
\[
A_f:P_1(U)\to P_0(V),\qquad
C_f:P_2(U)\to P_1(V),\qquad
B_h:P_0(U)\to P_2(V),
\]
where $f:U\to V$ is linear and $h:U\to E(V)$ is linear. The first two have respectively their sort-1 and sort-2 components equal to $f$. The last has sort-0 component $h$ and all others zero. Diagonal maps on $P_2$ have components $E(f)$ and $f$. There are no other nonzero off-diagonal Hom groups. Products of two successive off-diagonal types vanish:
\[
A_f C_g=0,\qquad B_h A_f=0,\qquad C_f B_h=0.
\]

Diagonal coefficient maps are tensored in the usual way on every displayed summand. Suppress the coefficient maps temporarily in order to give the off-diagonal formulas with second argument $P_0(W)$:
\begin{align*}
A\boxtimes P_0&=(0,q):P_1\longrightarrow P_0\oplus S,&
P_0\boxtimes A&=(A,-q):P_1\longrightarrow P_0\oplus S,\\
C\boxtimes P_0&=C:P_2\longrightarrow P_1,&
P_0\boxtimes C&=C:P_2\longrightarrow P_1.
\end{align*}
In each formula the parameter is the tensor product of the original parameters, and an $f$ is inserted as $f\otimes\id$ or $\id\otimes f$ in its indicated coefficient. For the third type put
\[
B_h\boxtimes P_0(W)=
\bigl(B_{\,t_{V,W}(h\otimes\id_W)},0\bigr):
P_0(U\otimes W)\oplus S(U\otimes W)\longrightarrow P_2(V\otimes W),
\]
and use the left-hand strength for $P_0(W)\boxtimes B_h$. Tensoring an off-diagonal map with $P_1$ or $P_2$ gives zero in every case. Extend these prescriptions additively to direct sums and then to retracts of free projectives.

Here is a full check that the prescriptions define a bifunctor. Naturality with diagonal coefficient maps is exactly naturality of $t$, or ordinary tensor naturality. For compositions of off-diagonal maps with the other argument $P_0$, the first variable gives respectively $qC=0$, $B(0,q)=0$, and $CB=0$; the second variable gives $(AC,-qC)=0$, $B(A,-q)=BA=0$, and $CB=0$. With a positive-index argument they vanish by the zero entries. For interchange of two off-diagonal maps, pairs $(A,A),(A,C),(C,A),(C,C),(B,B)$ have zero source, zero target, or a zero intermediate term in both orders. The pair $(B,A)$ and its reverse reduce to $BA=0$ after the formulas $(0,q)$ and $(A,-q)$; the pair $(B,C)$ and its reverse reduce to $CB=0$. These exhaust the nine ordered pairs. Thus interchange holds, including the pairs which failed in the two-sort attempted construction. All formulas preserve coproducts of the free projectives in each variable: this uses coproduct preservation by $P_i$ and $S$ and the canonical linear strength on coefficient maps.

\textbf{Right-exact extension.}
An additive functor on projectives extends to a right-exact functor on an abelian category with enough projectives by applying it to a projective presentation and taking the cokernel. Comparison maps and their homotopies give independence of the presentation and functoriality. The same construction in two variables extends the bifunctor just defined, using
\[
\coker\bigl((Q_1\boxtimes R_0)\oplus(Q_0\boxtimes R_1)
\longrightarrow Q_0\boxtimes R_0\bigr)
\]
for presentations $Q_1\to Q_0\to X\to0$ and $R_1\to R_0\to Y\to0$. Finite horseshoe presentations prove right exactness separately in each variable. This extension preserves arbitrary coproducts: coproducts of projective presentations remain right-exact presentations, coproducts of projectives are projective, and the prescribed functor on projectives preserves those coproducts. Exactness of coproducts at a monomorphism is not used. Hence the extended $\boxtimes$ is cocontinuous and $k$-bilinear in both variables.

Natural transformations between such right-exact functors are determined by their values on projectives and extend from compatible transformations there, again by presentations. We will use this fact in three or four variables to check all monoidal coherences.

\textbf{Associativity, symmetry and unit.}
The relevant finite-coordinate model is the category of cochain complexes in degrees $0,1,2$ over $k$, with the following tensor: retain the degree-zero tensor and terms having precisely one positive degree, and set every product of two positive-degree terms to zero. Thus degree $n>0$ is $(X_n\otimes Y_0)\oplus(X_0\otimes Y_n)$; the differential on degree zero is the pair of differentials, and the remaining differential acts on its sole positive factor. This is an associative symmetric tensor with unit $k$ in degree zero: a triple tensor retains exactly the terms with at most one positive factor. Its associativity and symmetry are the evident permutations of these terms. We use this coordinate model solely to specify and verify the finite matrices below; the new maps $B_h$ are checked separately.

For $P_0\boxtimes P_0=P_0\oplus S$, choose degree-one coordinates
\[
(x,y)=(p+s,p),
\]
where $x$ belongs to the first positive factor and $y$ to the second. Under these coordinates the symmetry is
\[
\begin{pmatrix}\id&i\\0&-\id\end{pmatrix}:
P_0(L)\oplus S(L)\longrightarrow P_0(L)\oplus S(L),
\]
followed by the coefficient swap $U\otimes V\to V\otimes U$. With one positive-index projective the symmetry is the coefficient swap on that projective; with two it is the unique map of zero objects.

From the presentation $P_2\to P_1\to S\to0$, right exactness gives
\[
S(U)\boxtimes P_0(V)=S(U\otimes V),\qquad
S(U)\boxtimes P_i(V)=0\quad(i>0),
\]
and the corresponding reversed formulas. Therefore a triple of $P_0$'s is $P_0(L)\oplus S(L)^{\oplus2}$ in either bracketing. Write the left coordinates as $(p,r,s)$, with $r$ the new $S$ from the outer tensor and $s$ the old $S$ from the first pair. Its three degree-one coordinates are
\[
(x,y,z)=(p+r+s,p+r,p).
\]
For the right bracketing they are $(p'+r',p'+s',p')$. Hence the associator is
\[
(p,r,s)\longmapsto(p,r+s,r).
\]
For a triple with exactly one positive index it is ordinary coefficient associativity on $P_i(L)$; with at least two positive indices both sides are zero.

These matrices are natural for the $A$, $C$ and diagonal maps by the finite-coordinate model. They are also natural for every $B_h$: those maps kill every $S$ summand, and the remaining equality is exactly the associativity or the symmetry-adjusted naturality of the canonical strength $t$. In the symmetry matrix the possible off-diagonal contribution is $B_h i=0$, since $A=i q$, $B_h A=0$, and $q$ is epimorphic (or directly because $\operatorname{Hom}(S(U),P_2(V))=0$). Thus no extra $E$-term is being discarded without verification.

The pentagon and hexagon can be checked on free projectives. With no positive index the displayed matrices arise from permutations and reassociations of the same finite-coordinate tensor, and therefore satisfy both diagrams. With exactly one positive index they are ordinary vector-space coefficient identities; with two or more positive indices every term is zero. Naturality for $B_h$ just checked permits extension from projectives, so these are coherences on all of $\mathcal B$.

Finally $1=\kappa k$ has presentation $P_1(k)\xrightarrow{A}P_0(k)\to1\to0$. Its left tensor with $P_0(V)$ is the cokernel of $(0,q)$, hence $P_0(V)$ by projection. Its right tensor is the cokernel of $(A,-q)$, hence $P_0(V)$ by $(\id,i)$. Tensoring $1$ with $P_i(V)$ for $i>0$ gives $P_i(V)$ directly. These unit identifications are natural for $B_h$ because $B_h i=0$ and $t_{V,k}=\id$; they are natural for the other generators by the same coordinate formulas. The triangle identities hold in the finite-coordinate model and extend by projective presentations. This completes the construction of a cocontinuous symmetric $k$-bilinear tensor with unit $\kappa k$. The exact-integer certificate in \texttt{code/C\_three\_sort\_coherence.py} checks the finite matrices, both hexagons, the pentagon, unit triangles, and generator relations in the finite model $E=\id$. That model is used only to make the extra $B$-map visible in a finite calculation; it does not verify $E(k)=0$ or the infinite-dimensional homology of this counterexample.

\textbf{The derived scalar calculation.}
The three displayed exact sequences give a projective resolution of $\kappa W$ with successive terms
\[
P_0(W),P_1(W),P_2(W),P_0(EW),P_1(EW),P_2(EW),\ldots.
\]
Applying $\phi$ leaves $E^n(W)$ in homological degree $3n$ and zero in the other degrees, with zero differentials. Thus
\[
L_{3n}\phi(\kappa W)=E^n(W),\qquad
L_{3n+1}\phi(\kappa W)=L_{3n+2}\phi(\kappa W)=0.
\]
The canonical degree-zero counit is $\id_W$. At $W=k$ this is the required normalization, while $L_3\phi(\kappa W)\ne0$ for infinite-dimensional $W$.

This category fails AB4. At $k$, the map $P_2(k)=S_2(k)\to P_1(k)$ is monic. Its coproduct indexed by a basis of an infinite-dimensional $W$ is $P_2(W)\to P_1(W)$, with nonzero kernel $\kappa E(W)$. Moreover $P_0(k)$ is not finitely presented: $\operatorname{Hom}(P_0(k),P_2(W))=E(W)$ whereas each $\operatorname{Hom}(P_0(k),P_2(k))$ is zero. If a set of finitely presented projective generators existed, its exact faithful evaluations would detect and prevent this failure of AB4. Thus the small-projective-generator corollary is not contradicted.


\textbf{The four manuscript hypotheses.}
Take $\mathcal M=\Ch_{\geq0}(\mathcal B)$, with the ordinary homological tensor induced by $\boxtimes$ and the usual signed symmetry. Only finite sums occur in each tensor degree. It is a complete and cocomplete abelian category with exact products and enough projectives; disk projectives give explicit degreewise epimorphisms onto every complex. The tensor is symmetric, $k$-bilinear and cocontinuous. Its unit is $\kappa k[0]$, its actual scalar-copower functor is $\kappa_{\mathcal M}(W)=\kappa W[0]$, and
\[
\varphi=\phi H_0:\mathcal M\longrightarrow\operatorname{Vect}_k
\]
is left adjoint to $\kappa_{\mathcal M}$. This scalar functor is exact and preserves all small limits, proving (a), and the disk construction proves (b). Projective resolutions of $\kappa W$ in $\mathcal B$, placed in internal degree zero, are projective resolutions of $\kappa W[0]$ in $\mathcal M$. Consequently the preceding calculation proves (c), and also
\[
L_3\varphi(\kappa_{\mathcal M}W)=E(W)\ne0.
\]

For (d) use $J=D^1(1)$ and its bottom inclusion $\alpha:1[0]\hookrightarrow J$. For every $X$, $\alpha\otimes X$ is degreewise split monic and $J\otimes X$ has the standard contraction $h(u_0\otimes x)=u_1\otimes x$, $h(u_1\otimes x)=0$, with $dh+hd=\id$. We verify total derived nullity without presuming exact coproducts. Any contractible nonnegative complex is the degreewise finite direct sum of disks $D^n(Z_{n-1})$, $n\geq1$. Choose projective resolutions of all the $Z_{n-1}$ in $\mathcal B$ and apply their disks. Their sum is a projective resolution in $\mathcal M$: at every internal degree it is a binary direct sum of exact resolutions. Every term of this resolution is a sum of disks and has $H_0=0$, so applying $\varphi$ gives the zero complex. Thus $L\varphi(J\otimes X)=0$ in all degrees, including zero. This proves (d) in its full strength.

\textbf{Failure of the asserted localization equivalence.}
For each vector space $W$, $\kappa_{\mathcal M}W$ is a $\mathcal W$-local object of $\mathcal M$: if $f$ is inverted by $L\varphi$, then $H_0L\varphi(f)=\varphi(f)$ is an isomorphism and the scalar adjunction makes $\operatorname{Hom}(f,\kappa_{\mathcal M}W)$ bijective. Therefore localization does not change maps into that object. For every $X\in\mathcal M$ we have naturally
\[
\operatorname{Hom}_{\mathcal M[\mathcal W^{-1}]}(X,\kappa_{\mathcal M}W)
=\operatorname{Hom}_k(\varphi X,W)
=\operatorname{Hom}_{D_{\geq0}(k)}(L\varphi X,W[0]).
\]
The last equality is the standard degree-zero truncation property, or follows directly from a nonnegative projective complex representing $L\varphi X$. If the asserted canonical functor were an equivalence, Yoneda, applied through that equivalence, would force its canonical map
\[
L\varphi(\kappa_{\mathcal M}W)\longrightarrow W[0]
\]
to be an isomorphism. Its third homology is $E(W)\ne0$ for infinite-dimensional $W$, a contradiction. If necessary we form the Gabriel--Zisman localization in a larger universe; a size failure of a smaller localization cannot restore the asserted equivalence.
\end{proof}


\paragraph{The necessary repair exposed by the counterexample.}
The last argument is general: if the canonical $L\varphi$-comparison is an equivalence and $\varphi\dashv\kappa$, then its canonical derived counit must satisfy
\[
L\varphi(\kappa V)\xrightarrow{\sim}V[0]
\quad\text{for every scalar module }V.
\]
Normalization at the single unit does not imply this requirement in the full manuscript scope. Claim~\ref{C26} fails precisely at that passage to infinite free coefficients. We do not assert that this necessary condition alone suffices; Claims~\ref{C11}, \ref{C13}, \ref{C15}, and~\ref{C23} retain their separately proved sufficient hypotheses.
